\chapter*{Abstract}
\addcontentsline{toc}{chapter}{Abstract}

\noindent\rule{\linewidth}{1pt}
\vspace{0.3cm}

\noindent The growing complexity of front-end applications developed as monoliths raises significant challenges in terms of maintainability, scalability, and team collaboration. Micro-frontend architecture addresses these challenges by decomposing a monolithic front-end into several autonomous sub-applications, each covering a distinct functional domain.

\medskip

\noindent This final-year project, carried out at Talan Tunisie, proposes an automated solution for migrating Angular monolithic applications toward a micro-frontend architecture using Module Federation. The solution takes the form of a multi-agent pipeline, where each specialized agent handles a precise step of the migration process: structural analysis of the source project, functional domain decomposition, micro-frontend architecture generation, component and service migration, progressive validation, and documentation production.

\medskip

\noindent The pipeline is orchestrated using LangGraph, which models the workflow as a stateful graph with controlled transitions between agents. Each intermediate artifact is validated before the next step begins, ensuring traceability and limiting error propagation throughout the process. The result is an executable, structured, and documented micro-frontend project ready for use by the development team.

\vspace{0.5cm}

\noindent\textbf{Keywords:} micro-frontend, Angular, Module Federation, multi-agent system, LangGraph, automated migration, monolith, software architecture.

\vspace{0.4cm}
\noindent\rule{\linewidth}{1pt}

\newpage

\section*{Résumé}
\addcontentsline{toc}{section}{Résumé}

\noindent\rule{\linewidth}{1pt}
\vspace{0.3cm}

\noindent La complexité croissante des applications front-end développées sous forme monolithique pose des défis importants en matière de maintenabilité, d'évolutivité et de collaboration entre équipes. L'architecture micro-frontend répond à ces enjeux en décomposant une application front-end en plusieurs sous-applications autonomes, chacune couvrant un domaine fonctionnel distinct.

\medskip

\noindent Ce projet de fin d'études, réalisé au sein de Talan Tunisie, propose une solution automatisée pour la migration d'applications Angular monolithiques vers une architecture micro-frontend basée sur Module Federation. La solution se présente sous la forme d'un pipeline multi-agents, où chaque agent spécialisé prend en charge une étape précise du processus : analyse structurelle du projet source, découpage fonctionnel, génération de l'architecture micro-frontend, migration des composants et services, validation progressive et production de documentation.

\medskip

\noindent Le pipeline est orchestré à l'aide de LangGraph, qui modélise le workflow sous forme d'un graphe à état avec des transitions contrôlées entre agents. Chaque artefact intermédiaire est validé avant le passage à l'étape suivante, garantissant la traçabilité et limitant la propagation des erreurs tout au long du processus. Le résultat est un projet micro-frontend exécutable, structuré et documenté, directement exploitable par l'équipe de développement.

\vspace{0.5cm}

\noindent\textbf{Mots clés :} micro-frontend, Angular, Module Federation, système multi-agents, LangGraph, migration automatisée, monolithe, architecture logicielle.

\vspace{0.4cm}
\noindent\rule{\linewidth}{1pt}
